\section{Verification, Acceptance, and Certification} \label{sec:verify-accept-certify}

This section defines the quality metrics for calibration products and ISR outputs.
Here we define the criteria used to promote products to ``accepted'' and ``certified''.

\subsection{Data products, Distribution, and User Access} \label{sec:verify-accept-certify:data-products}
Describe the output schema, distribution process (exporting/importing, rucio, TAXICAB tracking), and user access.

We enumerate the primary data products (calibration frames, ISR outputs, masks, and metadata tables), describe naming/versioning conventions, and summarize the provenance tracked for each product (inputs, software versions, and configuration).

Briefly describe how users discover and retrieve products (e.g., via a data butler/repository interface) and recommended query patterns.

Be sure to mention that header metadata contains information about whether a particular ISR step was applied/not applied and which calibration dataset was used.

\subsection{Verification} \label{sec:verify-accept-certify:verification}

We list (possibly in a table) the automated verification tests for each data product (sanity checks, analysis plots, detailed plots, distributional checks, and science test images) and the thresholds that determine acceptance/rejection of different calibration data products.

\subsection{Acceptance} \label{sec:verify-accept-certify:acceptance}
This section will describe the calibration acceptance process (TAXICAB).

Describe how verified products are compared against earlier baselines to reveal regressions or deviations, and how acceptance decisions are documented.

\subsection{Certification} \label{sec:verify-accept-certify:certification}
Describe the final certification step in the context of the data ``butler'' and file system storage.

\subsection{Data Monitoring} \label{sec:verify-accept-certify:data-monitoring}
Discuss how we monitor incoming streams of new data for quality (e.g.\ rubinTV, plot navigator, trend monitoring with chronograph).

\section{Limitations \& Known Issues} \label{sec:limitations}

The ISR pipeline developed for LSSTCam removes the dominant instrumental signatures, but several effects remain partially corrected, unmodeled, or unmasked.
The paragraphs below summarize known limitations (instrumental systematics not fully corrected, edge cases, validity constraints, etc.) for interpreting DP2 data, our recommendations for handling them in the context of DP2 science, and plans for near-term improvements expected before Data Release~1 (DR1).

\paragraph{Flares}
Some science images contain localized, bright flare-like structures.
They appear as elongated trails that begin at a detector edge and fade with a wavlength-dependent power-law profile consistent with attenuation of light in silicon at that wavelength.
The light is likely scattered from bright stars: focused starlight may reflect from the beveled edge of the silicon mount and enter a neighboring detector through its side, although this mechanism is not yet confirmed.
Similar structures appear in some flat-field images as many small edge flares (referred to as a ``comb" pattern), possibly linked to the rough-cut silicon edge.
So far, these features have been seen only in y-band, with both focused and unfocused illumination, and they are transient, consistent with a dependence on the chance alignment of bright stars with the focal-plane array.
They are not yet fully characterized or masked in DP2 processing.
Science users should treat anomalously bright detections/extended objects near detector edges with caution in y-band single-visit and difference images, and verify suspect sources by visual inspection.

\paragraph{Residual non-linearity wiggles}
The double-spline non-linearity model (\S\ref{sec:processing:calibration:linearity}) removes the dominant signal-dependent departure from linear response, but small quasi-oscillatory residuals remain as a function of signal level.
They appear as low-amplitude, signal-dependent structure in linearized images.
These ``wiggles'' are irregular, clearly observed beyond noise fluctuations, and they do not correlate with measured properties of the input flat-field images (e.g. time, temperature) or with channel on a detector or across the focal plane.
These residuals are well below current science requirements and are left uncorrected in DP2, but further characterization is in progress.

\paragraph{Differential non-linearity (DNL)}
The DP2 linearizer corrects average amplifier non-linearity but does not yet apply a full low-signal differential non-linearity correction.
This limitation reflects the difficulty of delivering flat-field illumination below 100 ADU/pixel with the in-dome calibration system.
Users should exercise caution when relying on photometric or astrometric measurements at low signal levels, in particular when relying on u-band sky-background estimates from later processing steps, where backgrounds are typically at or below 100 ADU/pixel.
Count dithering in ISR (\S\ref{sec:isr-science:overview}) provides interim mitigation until a dedicated DNL correction is measured and deployed.

\paragraph{Detector vignetting models}
Flat-field calibrations measure relative pixel and filter quantum efficiencies and throughputs under uniform illumination, but they require an accurate model of optical vignetting across the focal plane.
Improved vignetting models, informed by larger calibration datasets and cross-checks against on-sky stellar photometry, are planned for future calibration epochs to reduce focal-plane non-uniformity.

\paragraph{Color-dependent brighter-fatter corrections}
The brighter-fatter effect is color dependent because its amplitude depends on photon conversion depth in silicon and on the path length over which the transverse electric field acts on drifting electrons.
Laboratory and commissioning measurements show that the effect is chromatic \citep[$a_{00}^{y-band} = 0.94 \times a_{00}^{r-band}$]{broughton2024}.
The electrostatic brighter-fatter model used in DP2 (\S\ref{sec:processing:calibration:bf}) is fit from PTC covariances and applies a single wavelength-averaged correction per detector, derived in the band of the PTC input flats.
Because the effect is strongest near the bottom of the potential well and the front of the chip, calibrations for u through i bands are approximately consistent as most light in those bands converts to photoelectrons within the first few microns.
Separate calibrations are appropriate for z- and y-band light, which typically  convert far deeper into the silicon.
Boundary-shift calibrations can be color-corrected by assuming a flat incident spectral energy distribution in each filter band and a depth-conversion model in silicon.
These corrections were not enabled in DP2 because validation was incomplete, but they are expected for DR1.
Until then, PSF photometry and shape measurements of bright sources, especially in redder bands, may retain small flux- and filter-dependent biases.